\subsection{Assumptions}
    \textit{In this section, I would like to shortly describe the assumptions made by the Markowitz model. For example, that the investor is risk-averse. Some of the assumptions are economic assumptions that do not affect the result of the mathematical optimisation problem, but some others do. Moreover, it is a common practice in the financial industry to document Model assumptions and limitations in all model documentations.}
    
    \begin{definition}
        testest
    \end{definition}
    
    \begin{theorem}[]
        werwerew
    \end{theorem}
    
    References: \cite{markowitzorg}

\subsection{N-Assets Optimal Portfolio}\label{nassetssec}
    \textit{Here I will give a mathematical description of the optimisation problem presented by Markowitz, including some short proofs of the key points used to formulate the optimal solution as the solution of a linear equation system.}
    
    References: \cite{markowitzorg}, \cite{berndscherer}
    
\subsection{Practical Implementation}
    \textit{In this section, I will describe the optimal portfolios computed within the practical implementation framework}
    
    \subsubsection{Markowtiz's Curse}
        \textit{Using the results of the practical implementation, the concept of Markowitz's is defined here. Moreover, I will provide some mathematical explanation by computing the condition number of the covariance matrices used in the practical implementation and relating them to the error bounds we saw in Numerical Linear Algebra I}
        
        References: \cite{lopezmain}